\documentclass[11pt]{article}
\usepackage[margin=0.8in]{geometry}
\usepackage{amsmath}
\usepackage{amssymb}
\usepackage{enumitem}
\usepackage{fancyhdr}
\usepackage{multicol}

\pagestyle{fancy}
\fancyhf{}
\lhead{AP Statistics - Unit 2.9}
\rhead{Multiple Choice Practice}
\cfoot{\thepage}

\setlength{\parindent}{0pt}
\setlength{\parskip}{8pt}

\begin{document}

\begin{center}
{\Large \textbf{Topic 2.9: Analyzing Departures from Linearity}}\\[0.5em]
{\large Multiple Choice Questions}\\[1em]
Name: \underline{\hspace{3in}} \quad Period: \underline{\hspace{1in}}
\end{center}

\textbf{Instructions:} Select the best answer for each question. All questions have exactly one correct answer.

\begin{enumerate}

\item A point in a scatter plot has an $x$-value that is three standard deviations above the mean of all $x$-values in the dataset. This point is best classified as:
\begin{enumerate}[label=(\Alph*)]
\item A point with high correlation coefficient that affects the regression line minimally
\item A high leverage point that may substantially affect the regression line parameters
\item An outlier that primarily changes the coefficient of determination when removed
\item A residual point that indicates poor model fit but has minimal influence
\end{enumerate}

\item When a particular data point is removed from a linear regression analysis, the correlation coefficient $r$ increases from 0.65 to 0.89. This point is most likely:
\begin{enumerate}[label=(\Alph*)]
\item An outlier with a large residual that doesn't follow the general trend
\item A high leverage point with an unusual $x$-value following the trend
\item A low leverage point near the mean of both variables in the dataset
\item An influential point that primarily affects the $y$-intercept parameter only
\end{enumerate}

\item In examining a residual plot for a linear model, you observe a distinct curved pattern. The most appropriate next step would be to:
\begin{enumerate}[label=(\Alph*)]
\item Remove all points with residuals greater than two standard deviations
\item Transform one or both variables to achieve a more linear relationship
\item Increase the sample size to reduce the effect of individual outliers
\item Recalculate the regression line using a different statistical software package
\end{enumerate}

\item A data point has an $x$-value close to $\bar{x}$ but a residual of magnitude 4.2 when all other residuals are less than 1.5. This point should be classified as:
\begin{enumerate}[label=(\Alph*)]
\item A high leverage point that substantially affects the regression slope
\item An outlier that may significantly impact the correlation coefficient
\item A typical point that follows the established pattern in the data
\item An influential point that primarily changes the $x$-intercept value
\end{enumerate}

\item After applying a logarithmic transformation to the explanatory variable, the $R^2$ value increases from 0.42 to 0.78. This suggests that:
\begin{enumerate}[label=(\Alph*)]
\item The original relationship was already linear and transformation was unnecessary
\item The transformed model explains more variation in the response variable
\item The correlation coefficient decreased while the determination coefficient increased
\item The residual plot will show more pronounced patterns after transformation
\end{enumerate}

\item A point is identified as both a high leverage point and an outlier. Removing this point will most likely:
\begin{enumerate}[label=(\Alph*)]
\item Change only the correlation but leave slope and intercept unchanged
\item Affect the slope, $y$-intercept, and correlation coefficient substantially
\item Impact only the $y$-intercept while maintaining the original slope
\item Reduce the standard error without changing any regression parameters
\end{enumerate}

\item The mean point $(\bar{x}, \bar{y})$ in a dataset has which special property regarding the least-squares regression line?
\begin{enumerate}[label=(\Alph*)]
\item It always has the largest residual among all points in the dataset
\item It always lies exactly on the least-squares regression line equation
\item It determines the correlation coefficient but not the regression parameters
\item It represents the point with maximum leverage in the regression model
\end{enumerate}

\item Income data is often right-skewed because of a few high earners. When studying income versus another variable, applying a log transformation to income typically:
\begin{enumerate}[label=(\Alph*)]
\item Increases the skewness while preserving the original relationship pattern
\item Creates more outliers by emphasizing differences in high values
\item Reduces extreme values while maintaining the relative ordering of data
\item Eliminates all influential points from the regression analysis completely
\end{enumerate}

\item A researcher notices that removing points with $x$-values near $\bar{x}$ has minimal effect on the regression line. These points are best described as:
\begin{enumerate}[label=(\Alph*)]
\item High leverage points with substantial influence on model parameters
\item Low leverage points with minimal influence on regression coefficients  
\item Outliers that significantly affect the correlation coefficient value
\item Influential points that determine the strength of the relationship
\end{enumerate}

\item When assessing whether a transformation improved model fit, which combination of evidence would be most convincing?
\begin{enumerate}[label=(\Alph*)]
\item Decreased $R^2$ value and more random residual plot patterns
\item Increased $R^2$ value and more apparent randomness in residuals
\item Unchanged $R^2$ value but significantly different regression coefficients
\item Lower correlation coefficient with reduced variability in residuals
\end{enumerate}

\item A point with coordinates $(x_i, y_i)$ where $x_i$ is far from $\bar{x}$ but $y_i$ falls exactly on the regression line would be classified as:
\begin{enumerate}[label=(\Alph*)]
\item An outlier because it has an extreme value in the dataset
\item A high leverage point that follows the general trend perfectly
\item Both an outlier and high leverage point with maximum influence
\item Neither influential nor unusual since it fits the model well
\end{enumerate}

\item The residual plot for untransformed data shows a clear parabolic pattern. After log-transforming the $x$-variable, the residual plot shows random scatter. This indicates:
\begin{enumerate}[label=(\Alph*)]
\item The transformation created artificial patterns that weren't originally present
\item The linear model is now less appropriate for the transformed data
\item The transformation successfully linearized the relationship between variables
\item The original model was actually better despite the curved pattern
\end{enumerate}

\item An influential point is defined as a point that, when removed from the dataset:
\begin{enumerate}[label=(\Alph*)]
\item Always has both high leverage and a large residual value
\item Substantially changes slope, intercept, or correlation of the model
\item Must be located more than two standard deviations from the mean
\item Only affects the coefficient of determination but not other parameters
\end{enumerate}

\item Two models are fit to the same data: Model A uses untransformed variables with $R^2 = 0.55$, while Model B uses log-transformed $x$ with $R^2 = 0.82$. This suggests:
\begin{enumerate}[label=(\Alph*)]
\item Model A is superior because it uses the original data values
\item Model B better captures the relationship between the two variables
\item Both models are equally valid for making predictions about $y$
\item The transformation introduced bias that inflated the $R^2$ value
\end{enumerate}

\item A student claims that any point far from other points must be influential. This statement is:
\begin{enumerate}[label=(\Alph*)]
\item Correct, because distance from other points defines influence completely
\item Incorrect, because a high leverage point might still follow the trend
\item Correct, because isolated points always have large residual values
\item Incorrect, because influence depends only on the residual magnitude
\end{enumerate}

\item When examining the effect of country wealth on life expectancy, researchers found a curved relationship. The most appropriate transformation would likely be:
\begin{enumerate}[label=(\Alph*)]
\item Squaring the life expectancy values to emphasize larger differences
\item Taking the logarithm of wealth to reduce right skewness
\item Exponentiating both variables to create a linear relationship  
\item Using reciprocals of all values to reverse the relationship direction
\end{enumerate}

\item A data point has the largest residual in a dataset but an $x$-value exactly equal to $\bar{x}$. This point:
\begin{enumerate}[label=(\Alph*)]
\item Cannot be influential because it has low leverage by definition
\item Is likely an outlier that may affect the correlation coefficient
\item Must be a high leverage point due to its large residual
\item Will primarily change the slope but not the intercept value
\end{enumerate}

\item The primary purpose of transforming variables in regression analysis is to:
\begin{enumerate}[label=(\Alph*)]
\item Eliminate all outliers and influential points from the dataset
\item Increase the number of significant digits in correlation calculations
\item Create a more linear relationship suitable for linear regression methods
\item Ensure that all residuals have exactly the same magnitude
\end{enumerate}

\item A regression analysis yields $\hat{y} = 15 + 3x$ with $R^2 = 0.45$. After removing one point, the new equation is $\hat{y} = 8 + 3.1x$ with $R^2 = 0.44$. The removed point was most likely:
\begin{enumerate}[label=(\Alph*)]
\item A high leverage point that primarily affected the $y$-intercept
\item An outlier that substantially influenced the correlation coefficient
\item A low leverage point with minimal impact on model parameters
\item An influential point that changed both slope and correlation significantly
\end{enumerate}

\item In the context of regression analysis, which statement about the mean point $(\bar{x}, \bar{y})$ is always true?
\begin{enumerate}[label=(\Alph*)]
\item It has a residual of zero and maximum leverage value
\item It lies on the regression line with a residual of zero
\item It determines the slope but not the intercept of the line
\item It represents the most influential point in the dataset
\end{enumerate}

\item A dataset shows a strong positive curved relationship. After applying $\log(x)$ transformation, the residual plot still shows a pattern. The next step should be to:
\begin{enumerate}[label=(\Alph*)]
\item Accept that linear regression is inappropriate for this data
\item Try transforming the $y$-variable or both variables instead
\item Remove all high leverage points and refit the model
\item Conclude that the log transformation made the problem worse
\end{enumerate}

\item Points that are close to $\bar{x}$ on the $x$-axis are called low leverage points because they:
\begin{enumerate}[label=(\Alph*)]
\item Have minimal ability to influence the slope of the regression line
\item Always have small residuals and follow the general trend closely
\item Cannot be outliers regardless of their $y$-values in the dataset
\item Determine the correlation but not the regression coefficients
\end{enumerate}

\item A transformation is considered successful when comparing residual plots if:
\begin{enumerate}[label=(\Alph*)]
\item The transformed plot shows a clear linear pattern in residuals
\item The transformed plot displays more random scatter around zero
\item All residuals in the transformed plot have the same sign
\item The range of residuals increases substantially after transformation
\end{enumerate}

\item The distinction between an outlier and a high leverage point in regression is that:
\begin{enumerate}[label=(\Alph*)]
\item Outliers have unusual $y$-values while high leverage points have unusual $x$-values
\item Outliers have large residuals while high leverage points have extreme $x$-values
\item Outliers affect correlation while high leverage points affect only intercepts
\item Outliers are always influential while high leverage points are never influential
\end{enumerate}

\item When both squaring and logarithmic transformations are possible for linearizing data, the choice between them should be based on:
\begin{enumerate}[label=(\Alph*)]
\item The original correlation coefficient before any transformation is applied
\item Which transformation produces higher $R^2$ and better residual patterns
\item The number of outliers present in the original dataset
\item Whether the relationship is positive or negative in direction
\end{enumerate}

\end{enumerate}

\newpage
\section*{Answer Key}

\begin{multicols}{2}
\begin{enumerate}
\item \textbf{B} - A high leverage point has an unusual $x$-value (far from $\bar{x}$)
\item \textbf{A} - Large change in $r$ indicates an outlier with large residual
\item \textbf{B} - Transform variables when residual plot shows patterns
\item \textbf{B} - Close to $\bar{x}$ with large residual = outlier
\item \textbf{B} - Higher $R^2$ means more variation explained
\item \textbf{B} - Both outlier and high leverage affects all parameters
\item \textbf{B} - Mean point always lies on regression line
\item \textbf{C} - Log reduces extreme values, preserves order
\item \textbf{B} - Points near $\bar{x}$ are low leverage points
\item \textbf{B} - Increased $R^2$ + random residuals = better fit
\item \textbf{B} - High leverage but on the line (residual = 0)
\item \textbf{C} - Pattern removed = successful linearization
\item \textbf{B} - Influential = changes slope, intercept, or $r$
\item \textbf{B} - Higher $R^2$ indicates better model fit
\item \textbf{B} - High leverage can follow trend (not influential)
\item \textbf{B} - Log transformation for right-skewed income
\item \textbf{B} - At $\bar{x}$ with large residual = outlier
\item \textbf{C} - Transform to achieve linearity for linear methods
\item \textbf{A} - Large intercept change, small other changes
\item \textbf{B} - Mean point on line with zero residual
\item \textbf{B} - Try different transformations if pattern persists
\item \textbf{A} - Low leverage = minimal slope influence
\item \textbf{B} - Random scatter indicates good fit
\item \textbf{B} - Outliers: large residuals; high leverage: extreme $x$
\item \textbf{B} - Choose based on $R^2$ and residual patterns
\end{enumerate}
\end{multicols}

\end{document}